\section*{Chapter \ref{ch:b3:relativistic-dynamics} --- Relativistic Dynamics}

\begin{solution}{exo:b3:relativistic-dynamics:1}
Electron: $\gamma = 7.09$; $pc = \gamma\beta\,mc^2 = \qty{3.59}{MeV}$,
so $p = \qty{3.59}{MeV}/c$; $E = \qty{3.62}{MeV}$, $E_k =
\qty{3.11}{MeV}$. Proton: $E = 938 + 1000 = \qty{1938}{MeV}$, $\gamma
= 2.07$, $pc = \sqrt{E^2 - (mc^2)^2} = \qty{1696}{MeV}$, $\beta =
pc/E = 0.875$.
\end{solution}

\begin{solution}{exo:b3:relativistic-dynamics:2}
(a) $\num{e-3} \times \num{9e16} = \qty{9e13}{J}$. (b) $\num{6e13}/
\num{9e16} \approx \qty{0.7}{g}$ --- the bomb converted less than a
gram. (c) $L/c^2 = \qty{4.2e9}{kg/s}$; over \qty{e10}{yr}
($\num{3.2e17}{}\,\unit{s}$): $\num{1.3e27}{}\,\unit{kg}$, about
$0.07\%$ of the Sun. (d) $\num{5e4}/\num{9e16} \approx
\qty{0.6}{ng}$ --- real, and forever unweighable.
\end{solution}

\begin{solution}{exo:b3:relativistic-dynamics:3}
(a) $p = P/c$ per second: force $F = P/c = \qty{1.7e-11}{N}$ of
recoil. (b) Reflection doubles the transfer: $F = 2\Phi A/c =
2 \times 1400 \times 100/\num{3e8} \approx \qty{0.9}{mN}$. (c) $a
\approx \qty{9e-5}{m/s^2}$; after \qty{2.6e6}{s}: $v \approx
\qty{240}{m/s}$ --- slow to start, but the engine never runs dry. (d)
Sunlight's momentum (with the solar wind) pushes the dust
continuously outward: the tail streams away from the Sun, not behind
the comet.
\end{solution}

\begin{solution}{exo:b3:relativistic-dynamics:4}
(a) $m = E_0/c^2$: $\num{0.511e6} \times \num{1.6e-19}/\num{9e16} =
\qty{9.1e-31}{kg}$. (b) $\num{e9} \times \num{1.6e-19}/\num{3e8} =
\qty{5.3e-19}{kg.m/s}$. (c) $\sqrt s = \qty{4}{GeV}$ of
centre-of-momentum energy. (d) With $c = 1$, energy, momentum and
mass share one unit and every formula sheds its $c$'s; to return to
SI, reinsert the unique power of $c$ that fixes the dimensions.
\end{solution}

\begin{solution}{exo:b3:relativistic-dynamics:5}
(a) $2E_k/mc^2$: $1.96$, $3.9$, $17.6$, $59$ --- absurd beyond the
first. (b) $1 - (1 + E_k/mc^2)^{-2}$: $0.74$, $0.89$, $0.990$,
$0.9989$. (c) $8.4/0.99946c = \qty{28.0}{ns}$; Newton's $7.7c$ would
have read \qty{3.7}{ns}. (d) The calorimeter proved the electrons
truly carried the full \qty{15}{MeV} into the target: the energy was
all there, and still the speed had ceased to grow --- work now buys
momentum and energy, not velocity.
\end{solution}

\begin{solution}{exo:b3:relativistic-dynamics:6}
(a) $m_\pi c^2 = \sqrt{p^2c^2 + m_\mu^2c^4} + pc$; isolate the root
and square: $pc = (m_\pi^2 - m_\mu^2)c^4/2m_\pi c^2 =
(139.6^2 - 105.7^2)/(2 \times 139.6) = \qty{29.8}{MeV}$. (b) $E_\mu =
\sqrt{29.8^2 + 105.7^2} = \qty{109.8}{MeV}$: $E_k = \qty{4.1}{MeV}$,
$\beta = 29.8/109.8 = 0.27$. (c) $E_\nu = pc = \qty{29.8}{MeV}$. (d)
Boosting: $E_{\text{lab}} = \gamma(E^* \pm \beta p^*c)$ with $\beta
\approx 1$: between $50 \times (109.8 - 29.8) = \qty{4.0}{GeV}$ and
$50 \times (109.8 + 29.8) = \qty{7.0}{GeV}$.
\end{solution}

\begin{solution}{exo:b3:relativistic-dynamics:7}
(a) For two massless quanta, $M^2c^4 = 2E_1E_2(1 - \cos\theta) =
2 \times 200^2 \times \tfrac12$: $M = \qty{200}{MeV}/c^2$. (b)
Exactly $m_{\pi^0}$ --- \omterm{def:b3:relativistic-dynamics:four-momentum}{invariant mass} is the particle's mass, in
every frame. (c) Compute $M$ for every pair of photons in the event:
random pairs spread smoothly, true daughters of a particle pile up at
its mass --- the bump. (d) $\theta = 0$: $M = 0$; a parallel beam of
light has no rest frame --- it moves at $c$ as a whole, like each of
its photons.
\end{solution}

\begin{solution}{exo:b3:relativistic-dynamics:8}
(a) $\gamma = \num{6.8e12}/\num{9.38e8} = 7250$; $c - v \approx
c/2\gamma^2 = \qty{2.9}{m/s}$. (b) $f = v/L \approx \num{3e8}/27000 =
\qty{11.1}{kHz}$: eleven thousand laps per second. (c) $\num{3e14}
\times \qty{6.8}{TeV} = \qty{3.3e8}{J} \approx \qty{78}{kg}$ of TNT
--- in a hair-thin beam. (d) $2 \times \num{3.3e8}/\num{9e16} \approx
\qty{7}{\micro g}$: the working stock of matter-to-be.
\end{solution}

\begin{solution}{exo:b3:relativistic-dynamics:9}
(a) A photon's \omterm{def:b3:relativistic-dynamics:four-momentum}{invariant mass} is $0$; an $e^+e^-$ pair's is at least
$2m_{\text{e}}c^2$. The invariant cannot change in an isolated decay:
forbidden. (Equivalently: in no frame can momentum balance.) (b)
Head-on, $M^2c^4 = 4E\epsilon$; pair creation needs $Mc^2 \ge
2m_{\text{e}}c^2$, i.e.\ $E\epsilon \ge (m_{\text{e}}c^2)^2$. (c)
Another \qty{511}{keV} photon; for \qty{50}{keV}, a partner of
$0.511^2/0.05 = \qty{5.2}{MeV}$. (d) A TeV photon meets the
$(m_{\text{e}}c^2)^2/E \sim \qty{0.3}{eV}$ threshold on ordinary
starlight: the sky itself absorbs TeV $\gamma$-rays over cosmological
distances --- the universe is not transparent at all energies.
\end{solution}

\begin{solution}{exo:b3:relativistic-dynamics:10}
(a) $P_{e'} = P_\gamma + P_e - P_{\gamma'}$; square both sides
(invariants): $m_{\text{e}}^2c^4 = m_{\text{e}}^2c^4 + 2P_\gamma\!
\cdot\!P_e - 2P_{\gamma'}\!\cdot\!P_e - 2P_\gamma\!\cdot\!P_{\gamma'}$.
(b) With $P_\gamma\!\cdot\!P_e = E m_{\text{e}}$,
$P_{\gamma'}\!\cdot\!P_e = E'm_{\text{e}}$ and $P_\gamma\!\cdot\!
P_{\gamma'} = EE'(1 - \cos\theta)/c^2$:
$m_{\text{e}}c^2(E - E') = EE'(1 - \cos\theta)$, which in wavelengths
($E = hc/\lambda$) is the stated shift. (c) $h/m_{\text{e}}c =
\qty{2.43}{pm}$, at most \qty{4.9}{pm}: one part in $10^5$ of visible
light (invisible), several percent of a \qty{100}{pm} X-ray
(measured). (d) That a photon collides like a billiard ball ---
carrying momentum $h/\lambda$ exchanged in \emph{individual} events,
not merely energy in lumps.
\end{solution}

\begin{solution}{exo:b3:relativistic-dynamics:11}
(a) $M^2c^4 = m_{\text{p}}^2c^4 + 4E\epsilon$ head-on (the proton
ultrarelativistic); threshold at $M = m_\Delta$. (b) $E =
(1232^2 - 938^2)\,\unit{MeV^2}/(4 \times \num{6e-4}\,\unit{eV})
\approx \qty{2.7e20}{eV}$ for the mean photon --- the thermal tail of
the microwave background brings the effective cut-off to
$\sim\qty{6e19}{eV}$. (c) Protons above the cut-off lose energy to
the $\Delta$ within $\sim\qty{e8}{ly}$: the spectrum should end near
\qty{6e19}{eV} for distant sources --- the observed GZK suppression.
(d) Their sources must be both extraordinary accelerators and
cosmically \emph{nearby} --- within our supercluster --- which is
part of why their origin is still hunted.
\end{solution}

\begin{solution}{exo:b3:relativistic-dynamics:12}
(a) Energy: $m_{\text{i}}c^2 = \gamma m_{\text{f}}c^2 + E_\ell$;
momentum: $\gamma m_{\text{f}}\beta c = E_\ell/c$. Eliminate
$E_\ell$: $m_{\text{i}} = \gamma(1 + \beta)m_{\text{f}} =
m_{\text{f}}\sqrt{(1+\beta)/(1-\beta)}$. (b) $\sqrt{19} = 4.4$;
accelerating \emph{and} braking squares it: $19$. (c) The consumed
mass $m_{\text{i}} - m_{\text{f}} = 3.4\,m_{\text{f}}$ must be
annihilated fuel, half of it antimatter: \qty{1.7}{kg} of antimatter
per kilogram delivered (one way, no braking). (d) At nanograms per
year of world production, photon rockets remain arithmetic, not
engineering.
\end{solution}

\begin{solution}{pb:b3:relativistic-dynamics:1}
\textbf{1.} $E^2 - p^2c^2 = \gamma^2m^2c^4(1 - \beta^2) = m^2c^4$.
\textbf{2.} $E_{\text{tot}}$ and $\vect p_{\text{tot}}$ transform
like one particle's $E$ and $\vect p$ (they are sums), so the same
algebra gives an invariant; where $\vect p_{\text{tot}} = \vect 0$ it
is the total energy squared: $Mc^2 = E_{\text{CM}}$.
\textbf{3.} $(E + m_{\text{p}}c^2)^2 - p^2c^2 = m_{\text{p}}^2c^4 +
2Em_{\text{p}}c^2 + m_{\text{p}}^2c^4$.
\textbf{4.} $E \to m_{\text{p}}c^2$ gives $M = 2m_{\text{p}}$;
$E \gg m_{\text{p}}c^2$ gives $Mc^2 \to \sqrt{2Em_{\text{p}}c^2}$.
\textbf{5.} Any relative motion of the products adds kinetic energy
in the CM frame on top of their rest masses; the minimum $Mc^2$ ---
the threshold --- leaves them all at relative rest.
\textbf{6.} The lab momentum $\vect p_{\text{tot}} \neq \vect 0$ must
survive the collision: the products are condemned to move, and their
kinetic energy is locked away from mass-making.
\textbf{7.} $p + p \to p + p + p + \bar p$: charge $+2 \to +2$;
baryon number $+2 \to 1 + 1 + 1 - 1 = +2$. Making $\bar p$ alone, or
with anything less than a full extra baryon, breaks one of the two.
\textbf{8.} Four proton masses at relative rest: $Mc^2 =
4m_{\text{p}}c^2$.
\textbf{9.} $16m_{\text{p}}^2c^4 = 2m_{\text{p}}^2c^4 +
2Em_{\text{p}}c^2$: $E = 7m_{\text{p}}c^2$, $T = 6m_{\text{p}}c^2 =
\qty{5.63}{GeV}$.
\textbf{10.} New rest mass: $2m_{\text{p}}c^2$ out of $T =
6m_{\text{p}}c^2$ --- one third; the other two thirds fly on as the
products' kinetic energy, protected by momentum conservation.
\textbf{11.} Head-on, $Mc^2 = 2(T' + m_{\text{p}}c^2) =
4m_{\text{p}}c^2$: $T' = m_{\text{p}}c^2 = \qty{0.94}{GeV}$ per beam
--- six times less than the fixed-target beam, twelve times less
total kinetic energy.
\textbf{12.} Beams of the 1950s were far too tenuous to collide with
each other usefully; a solid target offers $10^{22}$ protons per
cubic centimetre. The price: $\qty{5.6}{GeV}$ instead of
$2 \times \qty{0.94}{GeV}$.
\textbf{13.} $E = 6.2 + 0.94 = \qty{7.14}{GeV}$; $pc =
\sqrt{7.14^2 - 0.938^2} = \qty{7.08}{GeV}$.
\textbf{14.} $p = \num{7.08e9} \times \num{1.6e-19}/\num{3e8} =
\qty{3.8e-18}{kg.m/s}$: $r = p/qB = \qty{15.1}{m}$ --- the machine as
built.
\textbf{15.} $\beta = pc/E = 0.991$; orbit $2\pi r = \qty{95}{m}$:
$f = \qty{3.1}{MHz}$.
\textbf{16.} At \qty{10}{MeV}, $\beta = 0.145$: $f =
\qty{0.46}{MHz}$. As the energy climbs, the speed (hence $f$) and the
momentum (hence the field needed at fixed radius) both change: field
and radio-frequency must ramp together --- the defining trick of the
synchrotron.
\textbf{17.} $\num{6.19e9}/1500 \approx \num{4e6}$ turns; at an
average megahertz-scale frequency, an acceleration cycle of the order
of one to two seconds.
\textbf{18.} The antiproton threshold $T = 6m_{\text{p}}c^2 =
\qty{5.6}{GeV}$, plus margin: the accountancy of this problem is
literally what the machine's energy --- and name --- were chosen for.
\textbf{19.} $r = p/qB$ contains no mass: the magnet bends equal
momenta equally, whatever the particle --- so the selected beam still
mixes pions, kaons and (rarely) antiprotons, all at
$\qty{1.19}{GeV}/c$.
\textbf{20.} $\bar p$: $E = \sqrt{1.19^2 + 0.938^2} =
\qty{1.51}{GeV}$, $\beta = 0.79$. $\pi^-$: $E = \sqrt{1.19^2 +
0.140^2} = \qty{1.20}{GeV}$, $\beta = 0.993$.
\textbf{21.} $t_{\bar p} = 12/(0.785c) = \qty{51}{ns}$; $t_\pi =
\qty{40}{ns}$: the \qty{11}{ns} gap demanded nanosecond timing ---
available, just, in 1955.
\textbf{22.} The counters were set to fire on the fast pions and stay
dark for the slow antiprotons (a velocity veto); with a $1$-in-$44000$
needle, only two independent signatures (time of flight \emph{and}
\v Cerenkov threshold) could exclude coincidence fakes.
\textbf{23.} Just above threshold the products are born almost at
relative rest: the reaction has almost no \omterm{def:b3:hamiltonian-mechanics:hamiltonian}{phase space}, while pion
production, far above \emph{its} threshold, is copious --- rarity was
guaranteed by the same kinematics that set the price.
\textbf{24.} $2m_{\text{p}}c^2 = \qty{1.9}{GeV}$ per annihilation,
typically into a handful of pions that decay onward to photons, muons
and neutrinos.
\textbf{25.} \omterm{def:b3:relativistic-dynamics:four-momentum}{Four-momentum} conservation priced the antiproton at
$6m_{\text{p}}c^2 = \qty{5.6}{GeV}$ on a fixed target; the
\qty{6.2}{GeV}, \qty{15.1}{m} Bevatron paid it; time of flight and a
\v Cerenkov veto found one antiworld particle per $44000$ impostors
--- and the 1959 Nobel Prize certified the ledger.
\end{solution}
